% GENERATED FILE. Edit canonical lessons, then run npm run generate:books.
% canonical-source-hash: fnv1a64:e7dcdb60cec25855
% canonical-lessons: MR-C208-aramkhurci, MR-C208-batan, MR-C208-khau, MR-C208-baksis, MR-C208-thamba

\chapter{Armchair, Button, Snack, Tip, Stop}
\label{ch:mr-armchair-button-snack-tip-stop}
\hlchaptermodality{208}

\begin{chapteropening}
\textbf{By the end of this chapter:} \emph{I can say an armchair, a button, a snack, a tip and a stop.}

Complete the last lesson of chapter 208: I can say an armchair, a button, a snack, a tip and a stop.
\end{chapteropening}

\section[\texorpdfstring{\mr{आरामखुर्ची} (ārāmkhurcī)}{आरामखुर्ची (ārāmkhurcī)}]{\mr{आरामखुर्ची} (ārāmkhurcī) — an armchair}

\label{lesson:MR-C208-aramkhurci}



\begin{quote}
Before the new one: say the Marathi for busy, then the Marathi for free.
\end{quote}

\subsection*{You'll want to know: \mr{आरामखुर्ची}}
\textbf{\mr{आरामखुर्ची}} — \emph{ārāmkhurcī} — \textquotedblleft{}an armchair\textquotedblright{}.

\begin{cousinweb}
\textbf{\mr{आरामखुर्ची}} (\emph{ārāmkhurcī}) is \textquotedblleft{}rest-chair\textquotedblright{}: \emph{ārām}, \textquotedblleft{}rest, ease\textquotedblright{}, usually traced to Persian \emph{ārām}, and \textbf{\mr{खुर्ची}}, \textquotedblleft{}chair\textquotedblright{}, from Arabic \emph{kursī}.
\end{cousinweb}

\begin{grammarlens}[title={What you've built}]
One more word for this chapter.
\end{grammarlens}

\subsection*{Your turn}
\begin{itemize}
\raggedright
  \item \emph{Say it:} \emph{ārāmkhurcī}
  \item \emph{Say it:} \emph{ārāmkhurcī}, once more
  \item \emph{Say it:} say \emph{ārāmkhurcī}
  \item \emph{Recall:} say the Marathi for busy, then the Marathi for free, then say \emph{ārāmkhurcī} again
\end{itemize}

\begin{tcolorbox}[breakable,colback=teal!4,colframe=teal!35!black,title={Before you move on}]
What does \mr{आरामखुर्ची} mean? (\textquotedblleft{}An armchair\textquotedblright{}.) Say it once more.
\end{tcolorbox}
\section[\texorpdfstring{\mr{बटण} (baṭaṇ)}{बटण (baṭaṇ)}]{\mr{बटण} (baṭaṇ) — a button}

\label{lesson:MR-C208-batan}



\begin{quote}
Before the new one: say the Marathi for free, then the Marathi for an armchair.
\end{quote}

\subsection*{You'll want to know: \mr{बटण}}
\textbf{\mr{बटण}} — \emph{baṭaṇ} — \textquotedblleft{}a button\textquotedblright{}.

\begin{cousinweb}
\textbf{\mr{बटण}} (\emph{baṭaṇ}) is a loan from English \emph{button}.
\end{cousinweb}

\begin{grammarlens}[title={What you've built}]
One more word for this chapter.
\end{grammarlens}

\subsection*{Your turn}
\begin{itemize}
\raggedright
  \item \emph{Say it:} \emph{baṭaṇ}
  \item \emph{Say it:} \emph{baṭaṇ}, once more
  \item \emph{Say it:} say \emph{baṭaṇ}
  \item \emph{Recall:} say the Marathi for free, then the Marathi for an armchair, then say \emph{baṭaṇ} again
\end{itemize}

\begin{tcolorbox}[breakable,colback=teal!4,colframe=teal!35!black,title={Before you move on}]
What does \mr{बटण} mean? (\textquotedblleft{}A button\textquotedblright{}.) Say it once more.
\end{tcolorbox}
\section[\texorpdfstring{\mr{खाऊ} (khāū)}{खाऊ (khāū)}]{\mr{खाऊ} (khāū) — a snack}

\label{lesson:MR-C208-khau}



\begin{quote}
Before the new one: say the Marathi for an armchair, then the Marathi for a button.
\end{quote}

\subsection*{You'll want to know: \mr{खाऊ}}
\textbf{\mr{खाऊ}} — \emph{khāū} — \textquotedblleft{}a snack\textquotedblright{}.

\begin{cousinweb}
\textbf{\mr{खाऊ}} (\emph{khāū}) is a tadbhava, a word inherited from Sanskrit and reshaped by sound change. It comes from \emph{khāduka}, \textquotedblleft{}biting\textquotedblright{} (CDIAL 3871). It is built on \emph{khād}, \textquotedblleft{}to eat\textquotedblright{}. CDIAL glosses the Marathi word \textquotedblleft{}eating\textquotedblright{}.
\end{cousinweb}

\begin{grammarlens}[title={What you've built}]
One more word for this chapter.
\end{grammarlens}

\subsection*{Your turn}
\begin{itemize}
\raggedright
  \item \emph{Say it:} \emph{khāū}
  \item \emph{Say it:} \emph{khāū}, once more
  \item \emph{Say it:} say \emph{khāū}
  \item \emph{Recall:} say the Marathi for an armchair, then the Marathi for a button, then say \emph{khāū} again
\end{itemize}

\begin{tcolorbox}[breakable,colback=teal!4,colframe=teal!35!black,title={Before you move on}]
What does \mr{खाऊ} mean? (\textquotedblleft{}A snack\textquotedblright{}.) Say it once more.
\end{tcolorbox}
\section[\texorpdfstring{\mr{बक्षीस} (bakṣīs)}{बक्षीस (bakṣīs)}]{\mr{बक्षीस} (bakṣīs) — a tip; a prize}

\label{lesson:MR-C208-baksis}



\begin{quote}
Before the new one: say the Marathi for a button, then the Marathi for a snack.
\end{quote}

\subsection*{You'll want to know: \mr{बक्षीस}}
\textbf{\mr{बक्षीस}} — \emph{bakṣīs} — \textquotedblleft{}a tip; a prize\textquotedblright{}.

\begin{cousinweb}
\textbf{\mr{बक्षीस}} (\emph{bakṣīs}) is a loan, usually traced to Persian \emph{bakhshish}, \textquotedblleft{}a gift\textquotedblright{}, from \emph{bakhshīdan}, \textquotedblleft{}to give\textquotedblright{}.
\end{cousinweb}

\begin{grammarlens}[title={What you've built}]
One more word for this chapter.
\end{grammarlens}

\subsection*{Your turn}
\begin{itemize}
\raggedright
  \item \emph{Say it:} \emph{bakṣīs}
  \item \emph{Say it:} \emph{bakṣīs}, once more
  \item \emph{Say it:} say \emph{bakṣīs}
  \item \emph{Recall:} say the Marathi for a button, then the Marathi for a snack, then say \emph{bakṣīs} again
\end{itemize}

\begin{tcolorbox}[breakable,colback=teal!4,colframe=teal!35!black,title={Before you move on}]
What does \mr{बक्षीस} mean? (\textquotedblleft{}A tip; a prize\textquotedblright{}.) Say it once more.
\end{tcolorbox}
\section[\texorpdfstring{\mr{थांबा} (thāmbā)}{थांबा (thāmbā)}]{\mr{थांबा} (thāmbā) — a stop}

\label{lesson:MR-C208-thamba}



\begin{quote}
Before the new one: say the Marathi for a snack, then the Marathi for a tip; a prize.
\end{quote}

\subsection*{You'll want to know: \mr{थांबा}}
\textbf{\mr{थांबा}} — \emph{thāmbā} — \textquotedblleft{}a stop\textquotedblright{}.

\begin{cousinweb}
\textbf{\mr{थांबा}} (\emph{thāmbā}) is the noun of \emph{thāṁbṇẽ}, \textquotedblleft{}to stop\textquotedblright{}, which comes from Sanskrit \emph{stambhate}, \textquotedblleft{}supports, arrests\textquotedblright{} (CDIAL 13683).
\end{cousinweb}

\begin{grammarlens}[title={What you've built}]
One more word for this chapter.
\end{grammarlens}

\subsection*{Your turn}
\begin{itemize}
\raggedright
  \item \emph{Say it:} \emph{thāmbā}
  \item \emph{Say it:} \emph{thāmbā}, once more
  \item \emph{Say it:} say \emph{thāmbā}
  \item \emph{Recall:} say the Marathi for a snack, then the Marathi for a tip; a prize, then say \emph{thāmbā} again
\end{itemize}

\begin{tcolorbox}[breakable,colback=teal!4,colframe=teal!35!black,title={Before you move on}]
What does \mr{थांबा} mean? (\textquotedblleft{}A stop\textquotedblright{}.) Say it once more.
\end{tcolorbox}
